\unitchapter{Paper 2 · Unit 3}{Programming Languages \& Computer Graphics}{p2-03}

\textbf{Expected questions: 8--10 · Target: 8+ · Type: C/C++ output prediction, OOP concepts, graphics algorithms \& transforms}

\begin{sylbox}

\begin{itemize}
\tightlist
\item[$\square$]
  Language design \& translation: paradigms, binding times, virtual computers, syntax, stages of translation
\item[$\square$]
  Elementary data types: scalar \& composite
\item[$\square$]
  Programming in C: tokens, data types, control, arrays, structures, unions, strings, pointers, functions, files, command-line args, preprocessor
\item[$\square$]
  OOP: class, object, inheritance, encapsulation, abstract class, polymorphism
\item[$\square$]
  C++: constructors/destructors, overloading, virtual functions, templates, exceptions, streams \& files
\item[$\square$]
  Web programming: HTML, DHTML, XML, scripting, Java, servlets, applets
\item[$\square$]
  Graphics: display devices, raster/random scan, line (DDA, Bresenham), circle \& ellipse (mid-point), polygon fill, boundary/flood fill
\item[$\square$]
  2-D transforms \& viewing: homogeneous coordinates, composite transforms, window-to-viewport, clipping (Cohen-Sutherland, Liang-Barsky, Sutherland-Hodgman)
\item[$\square$]
  3-D: polygon/quadric surfaces, splines, Bezier \& B-spline, illumination, shading, projections
\end{itemize}

\end{sylbox}

\needdiagram{figures/p2-03-00}{95pt}

\hypertarget{p2-03--1-language-concepts}{%
\section{1. Language Concepts}\label{p2-03--1-language-concepts}}

\needspace{100pt}

\hypertarget{p2-03--11-paradigms}{%
\subsection{1.1 Paradigms}\label{p2-03--11-paradigms}}

\diagramhere{figures/p2-03-00}

\begingroup\small

\begin{longtable}[]{@{}
  >{\raggedright\arraybackslash}p{(\columnwidth - 4\tabcolsep) * \real{0.1063}}
  >{\raggedright\arraybackslash}p{(\columnwidth - 4\tabcolsep) * \real{0.2777}}
  >{\raggedright\arraybackslash}p{(\columnwidth - 4\tabcolsep) * \real{0.2539}}@{}}
\toprule\noalign{}
\begin{minipage}[b]{\linewidth}\raggedright
\textbf{Language}
\end{minipage} & \begin{minipage}[b]{\linewidth}\raggedright
\textbf{Year/\allowbreak{}creator}
\end{minipage} & \begin{minipage}[b]{\linewidth}\raggedright
\textbf{Paradigm}
\end{minipage} \\
\midrule\noalign{}
\endhead
\bottomrule\noalign{}
\endlastfoot
FORTRAN & 1957, John Backus & First high-level, scientific \\
LISP & 1958, John McCarthy & Functional, AI \\
COBOL & 1959, Grace Hopper (CODASYL) & Business \\
ALGOL & 1958/60 & Block structure, BNF first used \\
Simula 67 & Dahl \& Nygaard & \textbf{First OO language} (classes) \\
Smalltalk & Alan Kay & Pure OO \\
C & 1972, Dennis Ritchie & Procedural \\
Prolog & 1972, Colmerauer & Logic \\
C++ & 1979--83, Bjarne Stroustrup & OO + procedural \\
Ada & 1980, US DoD & Concurrency (tasks, rendezvous) \\
Java & 1995, James Gosling (Sun) & OO, platform independent \\
Python & 1991, Guido van Rossum & Multi-paradigm \\
\end{longtable}

\endgroup

\needspace{132pt}

\hypertarget{p2-03--12-binding-time}{%
\subsection{1.2 Binding Time}\label{p2-03--12-binding-time}}

Binding = association of an attribute with an entity.

\begin{asciiblock}{42pt}
\begin{Verbatim}[fontsize=\fontsize{7.4}{9.0}\selectfont]
Language design time → Language implementation time → Compile time → Link time → Load time → Run time
  (meaning of *)          (size of int)                 (type of var)   (library code)  (address)    (value)
\end{Verbatim}
\end{asciiblock}

\begin{itemize}
\tightlist
\item
  \textbf{Static binding} (early) before run time; \textbf{dynamic binding} (late) at run time.
\item
  \textbf{Static scoping (lexical)}: binding by program text (C, Java). \textbf{Dynamic scoping}: by calling sequence (early LISP, Perl \texttt{local}).
\item
  Storage: static (globals), stack-dynamic (locals), explicit heap-dynamic (new/malloc), implicit heap-dynamic.
\end{itemize}

\needspace{135pt}

\hypertarget{p2-03--13-parameter-passing}{%
\subsection{1.3 Parameter Passing}\label{p2-03--13-parameter-passing}}

\begingroup\small

\begin{longtable}[]{@{}
  >{\raggedright\arraybackslash}p{(\columnwidth - 4\tabcolsep) * \real{0.3039}}
  >{\raggedright\arraybackslash}p{(\columnwidth - 4\tabcolsep) * \real{0.4466}}
  >{\raggedright\arraybackslash}p{(\columnwidth - 4\tabcolsep) * \real{0.1471}}@{}}
\toprule\noalign{}
\begin{minipage}[b]{\linewidth}\raggedright
\textbf{Method}
\end{minipage} & \begin{minipage}[b]{\linewidth}\raggedright
\textbf{Behaviour}
\end{minipage} & \begin{minipage}[b]{\linewidth}\raggedright
\textbf{Language}
\end{minipage} \\
\midrule\noalign{}
\endhead
\bottomrule\noalign{}
\endlastfoot
Call by value & Copy of value & C (default) \\
Call by reference & Address passed; changes visible & C++ \texttt{\&}, Pascal \texttt{var} \\
Call by value-result (copy-in/\allowbreak{}copy-out) & Copy in, copy back at return & Ada \texttt{in\ out} \\
Call by name & Textual substitution, evaluated each use (Jensen\textquotesingle s device) & ALGOL 60 \\
Call by need & Lazy evaluation, memoised & Haskell \\
\end{longtable}

\endgroup

\textbf{Classic question}:

\begin{asciiblock}{35pt}
\begin{Verbatim}[fontsize=\fontsize{8.8}{10.7}\selectfont]
void swap(int x, int y) { int t = x; x = y; y = t; }   // by value: caller unchanged
\end{Verbatim}
\end{asciiblock}

With \texttt{a\ =\ 1,\ i\ =\ 1,\ A{[}1{]}\ =\ 5}, call \texttt{f(i,\ A{[}i{]})} where f does \texttt{x\ =\ x\ +\ 1;\ y\ =\ y\ +\ 1;}:
by name → \texttt{A{[}2{]}} incremented (since i changed first); by reference → \texttt{A{[}1{]}} incremented.

\needspace{100pt}

\hypertarget{p2-03--14-data-types}{%
\subsection{1.4 Data Types}\label{p2-03--14-data-types}}

\begin{itemize}
\tightlist
\item
  \textbf{Scalar}: integer, real, character, Boolean, enumeration, pointer.
\item
  \textbf{Composite}: array, record/struct, union, string, set, list, file.
\item
  Strong typing (Java, Ada), weak typing (C). Type equivalence: \textbf{name} vs \textbf{structural}.
\item
  Type checking static (compile) vs dynamic (run time). Coercion = implicit type conversion.
\end{itemize}

\needspace{248pt}

\hypertarget{p2-03--2-c-programming--exam-traps}{%
\section{2. C Programming --- Exam Traps}\label{p2-03--2-c-programming--exam-traps}}

\needspace{208pt}

\hypertarget{p2-03--21-operator-precedence-high--low}{%
\subsection{2.1 Operator precedence (high → low)}\label{p2-03--21-operator-precedence-high--low}}

\begin{asciiblock}{153pt}
\begin{Verbatim}[fontsize=\fontsize{8.8}{10.7}\selectfont]
() [] -> .     ++ -- (postfix)
! ~ ++ -- + - * & (type) sizeof   (unary, right-to-left)
* / %
+ -
<< >>
< <= > >=
== !=
&   ^   |
&&  ||
?:              (right-to-left)
= += -= ...     (right-to-left)
,
\end{Verbatim}
\end{asciiblock}

\needspace{336pt}

\hypertarget{p2-03--22-output-prediction-examples}{%
\subsection{2.2 Output-prediction examples}\label{p2-03--22-output-prediction-examples}}

\begin{asciiblock}{281pt}
\begin{Verbatim}[fontsize=\fontsize{8.8}{10.7}\selectfont]
int i = 5;
printf("%d %d", i++, ++i);   // UNDEFINED behaviour (unsequenced modification)

int a = 10, b = 3;
printf("%d %d", a / b, a % b);   // 3 1
printf("%d", -7 % 3);            // -1 (sign follows dividend in C99)

int x = 0;
if (x = 5) printf("yes");        // prints yes (assignment, value 5)

char s[] = "NET";
printf("%zu %zu", sizeof(s), strlen(s));   // 4 3

int arr[] = {10, 20, 30};
int *p = arr;
printf("%d", *(p + 2));          // 30 ; arr[i] == *(arr + i) == i[arr]

int k = 3;
printf("%d", k << 2);            // 12 (multiply by 4)

static int counter() { static int c = 0; return ++c; }   // retains value: 1, 2, 3 ...

#define SQ(x) x*x
printf("%d", SQ(2+3));           // 2+3*2+3 = 11 (macro pitfall)
\end{Verbatim}
\end{asciiblock}

\needspace{144pt}

\hypertarget{p2-03--23-pointers}{%
\subsection{2.3 Pointers}\label{p2-03--23-pointers}}

\begin{asciiblock}{89pt}
\begin{Verbatim}[fontsize=\fontsize{8.8}{10.7}\selectfont]
int *p;          pointer to int
int **pp;        pointer to pointer
int *a[10];      array of 10 pointers to int
int (*a)[10];    pointer to an array of 10 ints
int *f();        function returning pointer to int
int (*f)();      pointer to function returning int
\end{Verbatim}
\end{asciiblock}

\begin{itemize}
\tightlist
\item
  Pointer arithmetic scales by \texttt{sizeof(type)}. \texttt{p\ +\ 1} on \texttt{int*} (4 bytes) adds 4.
\item
  \texttt{void*} generic pointer; dangling pointer (freed memory); NULL pointer; wild pointer (uninitialised).
\item
  Memory: \texttt{malloc} (uninitialised), \texttt{calloc} (zeroed, n × size), \texttt{realloc}, \texttt{free}.
\end{itemize}

\needspace{101pt}

\hypertarget{p2-03--24-structures-vs-unions}{%
\subsection{2.4 Structures vs Unions}\label{p2-03--24-structures-vs-unions}}

\begin{asciiblock}{46pt}
\begin{Verbatim}[fontsize=\fontsize{8.8}{10.7}\selectfont]
struct S { int i; char c; double d; };  // size = sum + padding (typically 16)
union  U { int i; char c; double d; };  // size = largest member (8) — members share memory
\end{Verbatim}
\end{asciiblock}

\needspace{135pt}

\hypertarget{p2-03--25-storage-classes}{%
\subsection{2.5 Storage Classes}\label{p2-03--25-storage-classes}}

\begingroup\small

\begin{longtable}[]{@{}
  >{\raggedright\arraybackslash}p{(\columnwidth - 8\tabcolsep) * \real{0.0753}}
  >{\raggedright\arraybackslash}p{(\columnwidth - 8\tabcolsep) * \real{0.0974}}
  >{\raggedright\arraybackslash}p{(\columnwidth - 8\tabcolsep) * \real{0.1273}}
  >{\raggedright\arraybackslash}p{(\columnwidth - 8\tabcolsep) * \real{0.0813}}
  >{\raggedright\arraybackslash}p{(\columnwidth - 8\tabcolsep) * \real{0.1753}}@{}}
\toprule\noalign{}
\begin{minipage}[b]{\linewidth}\raggedright
\textbf{Class}
\end{minipage} & \begin{minipage}[b]{\linewidth}\raggedright
\textbf{Scope}
\end{minipage} & \begin{minipage}[b]{\linewidth}\raggedright
\textbf{Lifetime}
\end{minipage} & \begin{minipage}[b]{\linewidth}\raggedright
\textbf{Default}
\end{minipage} & \begin{minipage}[b]{\linewidth}\raggedright
\textbf{Storage}
\end{minipage} \\
\midrule\noalign{}
\endhead
\bottomrule\noalign{}
\endlastfoot
auto & Block & Block & Garbage & Stack \\
register & Block & Block & Garbage & CPU register (no \texttt{\&}) \\
static & Block/file & Whole program & 0 & Data segment \\
extern & Global & Whole program & 0 & Data segment \\
\end{longtable}

\endgroup

\needspace{100pt}

\hypertarget{p2-03--26-files-command-line-preprocessor}{%
\subsection{2.6 Files, Command line, Preprocessor}\label{p2-03--26-files-command-line-preprocessor}}

\begin{itemize}
\tightlist
\item
  \texttt{FILE\ *fp\ =\ fopen("a.txt",\ "r")}; modes r, w, a, r+, w+, a+, rb. \texttt{fgetc}, \texttt{fputc}, \texttt{fgets}, \texttt{fprintf}, \texttt{fscanf}, \texttt{fread}, \texttt{fwrite}, \texttt{fseek}, \texttt{ftell}, \texttt{rewind}.
\item
  \texttt{int\ main(int\ argc,\ char\ *argv{[}{]})} --- \texttt{argv{[}0{]}} is the program name; \texttt{argc\ ≥\ 1}.
\item
  Preprocessor: \texttt{\#include}, \texttt{\#define}, \texttt{\#ifdef}, \texttt{\#ifndef}, \texttt{\#if}, \texttt{\#pragma}, \texttt{\#undef}; runs \textbf{before compilation}.
\end{itemize}

\needdiagram{figures/p2-03-01}{55pt}

\hypertarget{p2-03--3-oop-concepts}{%
\section{3. OOP Concepts}\label{p2-03--3-oop-concepts}}

\diagramhere{figures/p2-03-01}

\begin{asciiblock}{99pt}
\begin{Verbatim}[fontsize=\fontsize{8.8}{10.7}\selectfont]
 Inheritance types
 Single      Multiple      Multilevel     Hierarchical      Hybrid (diamond)
   A         A     B          A               A                  A
   │          ╲   ╱           │             ╱ │ ╲               ╱ ╲
   B            C             B            B  C  D             B   C
                              │                                 ╲ ╱
                              C                                  D
\end{Verbatim}
\end{asciiblock}

\begin{itemize}
\tightlist
\item
  Diamond problem in C++ solved by \textbf{virtual base class}. Java avoids multiple class inheritance (uses interfaces).
\end{itemize}

\needspace{135pt}

\hypertarget{p2-03--4-c-specifics}{%
\section{4. C++ Specifics}\label{p2-03--4-c-specifics}}

\begingroup\small

\begin{longtable}[]{@{}
  >{\raggedright\arraybackslash}p{(\columnwidth - 2\tabcolsep) * \real{0.1660}}
  >{\raggedright\arraybackslash}p{(\columnwidth - 2\tabcolsep) * \real{0.8340}}@{}}
\toprule\noalign{}
\begin{minipage}[b]{\linewidth}\raggedright
\textbf{Topic}
\end{minipage} & \begin{minipage}[b]{\linewidth}\raggedright
\textbf{Key point}
\end{minipage} \\
\midrule\noalign{}
\endhead
\bottomrule\noalign{}
\endlastfoot
Constructor & Same name as class, no return type; default, parameterised, \textbf{copy constructor} \texttt{X(const\ X\&)} \\
Destructor & \texttt{\textasciitilde{}X()}, no args, cannot be overloaded; called in reverse order of construction \\
Order of construction & Base → member objects → derived; destruction reverse \\
\texttt{this} pointer & Hidden pointer to invoking object; not available in static functions \\
Friend function & Not a member but accesses private data; not inherited \\
Static member & Shared by all objects; static function accesses only static members \\
Inline function & Expanded at call site \\
Function overloading & Same name, different parameter lists (return type alone not enough) \\
Operator overloading & Cannot overload \texttt{::}, \texttt{.}, \texttt{.*}, \texttt{?:}, \texttt{sizeof} \\
Virtual function & Run-time polymorphism via \textbf{vtable/vptr}; call through base pointer/\allowbreak{}reference \\
Pure virtual & \texttt{virtual\ void\ f()\ =\ 0;} → class becomes \textbf{abstract} (cannot be instantiated) \\
Virtual destructor & Needed when deleting derived object via base pointer \\
Templates & Generic programming: function \& class templates (compile-time polymorphism) \\
Exceptions & \texttt{try}, \texttt{throw}, \texttt{catch(...)} catches all \\
Streams & \texttt{cin} (istream), \texttt{cout} (ostream), \texttt{cerr} (unbuffered), \texttt{clog}; \texttt{ifstream}, \texttt{ofstream}, \texttt{fstream} \\
Access in inheritance & public base: public→\allowbreak{}public, protected→\allowbreak{}protected; protected base: both→\allowbreak{}protected; private base: both→\allowbreak{}private; private members never accessible \\
\end{longtable}

\endgroup

\begin{asciiblock}{67pt}
\begin{Verbatim}[fontsize=\fontsize{8.8}{10.7}\selectfont]
class Base { public: virtual void show() { cout << "Base"; } };
class Der : public Base { public: void show() { cout << "Derived"; } };
Base *b = new Der();
b->show();    // "Derived" (virtual)  — without virtual: "Base"
\end{Verbatim}
\end{asciiblock}

\needspace{135pt}

\hypertarget{p2-03--5-web-programming}{%
\section{5. Web Programming}\label{p2-03--5-web-programming}}

\begingroup\small

\begin{longtable}[]{@{}
  >{\raggedright\arraybackslash}p{(\columnwidth - 2\tabcolsep) * \real{0.1274}}
  >{\raggedright\arraybackslash}p{(\columnwidth - 2\tabcolsep) * \real{0.8726}}@{}}
\toprule\noalign{}
\begin{minipage}[b]{\linewidth}\raggedright
\textbf{Technology}
\end{minipage} & \begin{minipage}[b]{\linewidth}\raggedright
\textbf{Key facts}
\end{minipage} \\
\midrule\noalign{}
\endhead
\bottomrule\noalign{}
\endlastfoot
HTML & Markup; tags \texttt{\textless{}a\ href\textgreater{}}, \texttt{\textless{}img\ src\ alt\textgreater{}}, \texttt{\textless{}table\textgreater{}\textless{}tr\textgreater{}\textless{}td\textgreater{}}, \texttt{\textless{}form\ action\ method\textgreater{}}; HTML5 adds \texttt{\textless{}canvas\textgreater{}}, \texttt{\textless{}video\textgreater{}}, \texttt{\textless{}audio\textgreater{}}, \texttt{\textless{}section\textgreater{}}, \texttt{\textless{}article\textgreater{}} \\
DHTML & HTML + CSS + JavaScript + DOM → dynamic pages \\
CSS & Inline \textgreater{} internal \textgreater{} external (priority); selectors \texttt{.class}, \texttt{\#id} \\
XML & User-defined tags, data description, must be well-formed (single root, closed tags, case-sensitive); validated by \textbf{DTD/XSD}; XSLT for transformation, XPath for navigation \\
JavaScript & Client-side scripting; \texttt{document.getElementById} \\
Server-side & PHP, JSP, ASP.NET, Servlets \\
Java Applet & Runs in browser; life cycle: \texttt{init()\ →\ start()\ →\ paint()\ →\ stop()\ →\ destroy()}; no \texttt{main()} (now obsolete) \\
Servlet & Runs on server; life cycle: \texttt{init()\ →\ service()\ (doGet/doPost)\ →\ destroy()}; container e.g. Tomcat \\
Java & Bytecode on JVM ("write once run anywhere"); garbage collection; no pointers; \texttt{final}, \texttt{abstract}, \texttt{interface} \\
HTTP methods & GET (in URL, idempotent), POST (in body), PUT, DELETE \\
\end{longtable}

\endgroup

\needspace{175pt}

\hypertarget{p2-03--6-computer-graphics}{%
\section{6. Computer Graphics}\label{p2-03--6-computer-graphics}}

\needspace{135pt}

\hypertarget{p2-03--61-display-devices}{%
\subsection{6.1 Display Devices}\label{p2-03--61-display-devices}}

\begingroup\small

\begin{longtable}[]{@{}
  >{\raggedright\arraybackslash}p{(\columnwidth - 2\tabcolsep) * \real{0.2986}}
  >{\raggedright\arraybackslash}p{(\columnwidth - 2\tabcolsep) * \real{0.3019}}@{}}
\toprule\noalign{}
\begin{minipage}[b]{\linewidth}\raggedright
\textbf{Raster scan}
\end{minipage} & \begin{minipage}[b]{\linewidth}\raggedright
\textbf{Random (vector) scan}
\end{minipage} \\
\midrule\noalign{}
\endhead
\bottomrule\noalign{}
\endlastfoot
Beam sweeps every row (top to bottom) & Beam draws only the lines of the picture \\
Uses frame buffer (refresh buffer) & Uses display file (display list) \\
Realistic images, shading & Line drawings only, smooth lines \\
Aliasing (jaggies) & No aliasing \\
TV, monitors & Pen plotters, early CAD \\
\end{longtable}

\endgroup

\begin{asciiblock}{78pt}
\begin{Verbatim}[fontsize=\fontsize{8.8}{10.7}\selectfont]
Frame buffer size = Resolution × bits per pixel / 8   bytes
e.g. 1024 × 768, 24-bit colour → 1024×768×3 = 2.25 MB
Number of colours = 2^(bits per pixel)
Refresh: interlaced (odd/even lines alternately) vs non-interlaced
Aspect ratio = width : height
\end{Verbatim}
\end{asciiblock}

\begin{itemize}
\tightlist
\item
  CRT parts: electron gun, focusing \& deflection system, phosphor coated screen. \textbf{Persistence} = time phosphor glows. Shadow-mask (RGB) and beam-penetration (limited colours) colour CRTs.
\item
  LCD, LED, plasma (flat panel); DVST (Direct View Storage Tube). Colour lookup table saves memory.
\end{itemize}

\needspace{157pt}

\hypertarget{p2-03--62-line-drawing}{%
\subsection{6.2 Line Drawing}\label{p2-03--62-line-drawing}}

\textbf{DDA (Digital Differential Analyzer)}

\begin{asciiblock}{67pt}
\begin{Verbatim}[fontsize=\fontsize{8.8}{10.7}\selectfont]
dx = x2 − x1, dy = y2 − y1, steps = max(|dx|, |dy|)
xinc = dx/steps, yinc = dy/steps
repeat steps times: x += xinc, y += yinc, plot(round(x), round(y))
→ uses floating point and rounding (slower)
\end{Verbatim}
\end{asciiblock}

\textbf{Bresenham (\textbar m\textbar{} \textless{} 1)} --- integer only:

\begin{asciiblock}{57pt}
\begin{Verbatim}[fontsize=\fontsize{8.8}{10.7}\selectfont]
p0 = 2dy − dx
if pk < 0:  next = (xk+1, yk),     pk+1 = pk + 2dy
else:       next = (xk+1, yk+1),   pk+1 = pk + 2dy − 2dx
\end{Verbatim}
\end{asciiblock}

Example (20,10)→(30,18): dx=10, dy=8, p0 = 6, 2dy = 16, 2dy−2dx = −4.

\begin{asciiblock}{110pt}
\begin{Verbatim}[fontsize=\fontsize{8.8}{10.7}\selectfont]
 k  pk   plot
 0   6   (21,11)
 1   2   (22,12)
 2  −2   (23,12)
 3  14   (24,13)
 4  10   (25,14)
 5   6   (26,15)
 ...
\end{Verbatim}
\end{asciiblock}

\needspace{100pt}

\hypertarget{p2-03--63-mid-point-circle}{%
\subsection{6.3 Mid-point Circle}\label{p2-03--63-mid-point-circle}}

\begin{itemize}
\tightlist
\item
  8-way symmetry: compute one octant (x from 0 to x = y), reflect to 8.
\end{itemize}

\begin{asciiblock}{57pt}
\begin{Verbatim}[fontsize=\fontsize{8.8}{10.7}\selectfont]
p0 = 1 − r     (5/4 − r)
if pk < 0:  (xk+1, yk),     pk+1 = pk + 2xk+1 + 1
else:       (xk+1, yk−1),   pk+1 = pk + 2xk+1 + 1 − 2yk+1
\end{Verbatim}
\end{asciiblock}

Ellipse: \textbf{4-way symmetry}, two regions (slope −1 boundary).

\needspace{100pt}

\hypertarget{p2-03--64-filling}{%
\subsection{6.4 Filling}\label{p2-03--64-filling}}

\begin{itemize}
\tightlist
\item
  \textbf{Scan-line polygon fill}: intersect each scan line with edges, sort x, fill pairs; uses edge table \& active edge table.
\item
  \textbf{Boundary fill}: fills until boundary colour found. \textbf{Flood fill}: replaces a specific interior colour.
\item
  4-connected vs 8-connected (8-connected may leak through diagonal gaps).
\item
  Inside test: \textbf{odd-even (even-odd) rule}, \textbf{non-zero winding number rule}.
\end{itemize}

\needspace{161pt}

\hypertarget{p2-03--65-2-d-transformations-homogeneous-coordinates}{%
\subsection{6.5 2-D Transformations (Homogeneous coordinates)}\label{p2-03--65-2-d-transformations-homogeneous-coordinates}}

\begin{asciiblock}{106pt}
\begin{Verbatim}[fontsize=\fontsize{8.4}{10.2}\selectfont]
Translation        Scaling            Rotation (anticlockwise θ)
| 1  0  tx |       | sx 0  0 |        | cosθ  −sinθ  0 |
| 0  1  ty |       | 0  sy 0 |        | sinθ   cosθ  0 |
| 0  0  1  |       | 0  0  1 |        |  0      0    1 |

Reflection about x-axis: diag(1, −1, 1)      about y-axis: diag(−1, 1, 1)
about origin: diag(−1, −1, 1)                about y = x: swap x and y  [[0 1 0][1 0 0][0 0 1]]
X-shear: x' = x + shx·y                      Y-shear: y' = y + shy·x
\end{Verbatim}
\end{asciiblock}

\begin{itemize}
\tightlist
\item
  Homogeneous coordinates allow \textbf{translation as matrix multiplication} → composite transforms.
\item
  Rotation about arbitrary point (xr, yr): \textbf{T(xr, yr) · R(θ) · T(−xr, −yr)} (applied right-to-left).
\item
  Successive translations add; successive rotations add angles; successive scalings multiply.
\item
  Matrix multiplication is generally \textbf{not commutative} (except e.g. two rotations, two scalings, rotation \& uniform scaling).
\end{itemize}

\needdiagram{figures/p2-03-02}{55pt}

\hypertarget{p2-03--66-viewing--window-to-viewport}{%
\subsection{6.6 Viewing \& Window-to-Viewport}\label{p2-03--66-viewing--window-to-viewport}}

\diagramhere{figures/p2-03-02}

\begin{asciiblock}{46pt}
\begin{Verbatim}[fontsize=\fontsize{8.8}{10.7}\selectfont]
xv = xvmin + (xw − xwmin) · sx      sx = (xvmax − xvmin)/(xwmax − xwmin)
yv = yvmin + (yw − ywmin) · sy      sy = (yvmax − yvmin)/(ywmax − ywmin)
\end{Verbatim}
\end{asciiblock}

\needspace{168pt}

\hypertarget{p2-03--67-clipping}{%
\subsection{6.7 Clipping}\label{p2-03--67-clipping}}

\textbf{Cohen--Sutherland line clipping} --- 4-bit region code \textbf{TBRL} (Top, Bottom, Right, Left):

\begin{asciiblock}{78pt}
\begin{Verbatim}[fontsize=\fontsize{8.8}{10.7}\selectfont]
        1001 │ 1000 │ 1010
       ──────┼──────┼──────
        0001 │ 0000 │ 0010        window = 0000
       ──────┼──────┼──────
        0101 │ 0100 │ 0110
\end{Verbatim}
\end{asciiblock}

\begin{itemize}
\tightlist
\item
  Both codes 0000 → trivially accept. (code1 AND code2) ≠ 0 → trivially reject. Else clip against an edge and repeat.
\end{itemize}

\begingroup\small

\begin{longtable}[]{@{}
  >{\raggedright\arraybackslash}p{(\columnwidth - 2\tabcolsep) * \real{0.1741}}
  >{\raggedright\arraybackslash}p{(\columnwidth - 2\tabcolsep) * \real{0.3566}}@{}}
\toprule\noalign{}
\begin{minipage}[b]{\linewidth}\raggedright
\textbf{Algorithm}
\end{minipage} & \begin{minipage}[b]{\linewidth}\raggedright
\textbf{For}
\end{minipage} \\
\midrule\noalign{}
\endhead
\bottomrule\noalign{}
\endlastfoot
Cohen--Sutherland & Lines (region codes) \\
\textbf{Liang--Barsky} & Lines (parametric, faster) \\
Cyrus--Beck & Lines vs convex polygon window \\
Nicholl--Lee--Nicholl & Lines (fewest comparisons) \\
Mid-point subdivision & Lines (binary search) \\
\textbf{Sutherland--Hodgman} & Polygon clipping (convex window; edge by edge) \\
Weiler--Atherton & Polygon clipping (concave too) \\
\end{longtable}

\endgroup

\needspace{100pt}

\hypertarget{p2-03--68-3-d-graphics}{%
\subsection{6.8 3-D Graphics}\label{p2-03--68-3-d-graphics}}

\begin{itemize}
\tightlist
\item
  \textbf{Polygon surfaces}: vertex, edge, polygon tables; plane equation Ax + By + Cz + D = 0.
\item
  \textbf{Quadric surfaces}: sphere, ellipsoid, torus.
\item
  \textbf{Splines}: interpolation (passes through control points) vs approximation (near them).
\end{itemize}

\begingroup\small

\begin{longtable}[]{@{}
  >{\raggedright\arraybackslash}p{(\columnwidth - 2\tabcolsep) * \real{0.4024}}
  >{\raggedright\arraybackslash}p{(\columnwidth - 2\tabcolsep) * \real{0.3750}}@{}}
\toprule\noalign{}
\begin{minipage}[b]{\linewidth}\raggedright
\textbf{Bezier curve}
\end{minipage} & \begin{minipage}[b]{\linewidth}\raggedright
\textbf{B-spline curve}
\end{minipage} \\
\midrule\noalign{}
\endhead
\bottomrule\noalign{}
\endlastfoot
Degree = n − 1 for n control points & Degree independent of number of points \\
Passes through first \& last control points & Generally doesn\textquotesingle t pass through endpoints (uniform) \\
\textbf{Global control} --- moving one point changes whole curve & \textbf{Local control} \\
Lies within convex hull of control points & Lies within convex hull \\
Bernstein basis polynomials & B-spline basis functions \\
\end{longtable}

\endgroup

Cubic Bezier: P(t) = (1−t)³P₀ + 3t(1−t)²P₁ + 3t²(1−t)P₂ + t³P₃, 0 ≤ t ≤ 1.

\needspace{101pt}

\hypertarget{p2-03--69-illumination--shading}{%
\subsection{6.9 Illumination \& Shading}\label{p2-03--69-illumination--shading}}

\begin{asciiblock}{46pt}
\begin{Verbatim}[fontsize=\fontsize{8.8}{10.7}\selectfont]
I = Ia·ka  +  Il·kd·(N·L)  +  Il·ks·(R·V)ⁿ
    ambient     diffuse (Lambert)   specular (Phong, n = shininess)
\end{Verbatim}
\end{asciiblock}

\begingroup\small

\begin{longtable}[]{@{}
  >{\raggedright\arraybackslash}p{(\columnwidth - 4\tabcolsep) * \real{0.1282}}
  >{\raggedright\arraybackslash}p{(\columnwidth - 4\tabcolsep) * \real{0.2655}}
  >{\raggedright\arraybackslash}p{(\columnwidth - 4\tabcolsep) * \real{0.3893}}@{}}
\toprule\noalign{}
\begin{minipage}[b]{\linewidth}\raggedright
\textbf{Shading}
\end{minipage} & \begin{minipage}[b]{\linewidth}\raggedright
\textbf{Method}
\end{minipage} & \begin{minipage}[b]{\linewidth}\raggedright
\textbf{Quality}
\end{minipage} \\
\midrule\noalign{}
\endhead
\bottomrule\noalign{}
\endlastfoot
Flat (constant) & One intensity per polygon & Faceted \\
\textbf{Gouraud} & Interpolate \textbf{intensities} at vertices & Smooth, may miss specular highlights, Mach bands \\
\textbf{Phong} & Interpolate \textbf{normals} & Best quality, costlier \\
\end{longtable}

\endgroup

\needdiagram{figures/p2-03-03}{55pt}

\hypertarget{p2-03--610-projections}{%
\subsection{6.10 Projections}\label{p2-03--610-projections}}

\diagramhere{figures/p2-03-03}

\begin{itemize}
\tightlist
\item
  Perspective: realistic, foreshortening, parallel lines meet at \textbf{vanishing points}; does not preserve parallelism/size.
\item
  Cavalier (receding lines full length, 45°) vs Cabinet (half length, 63.4°).
\item
  Hidden surface removal: \textbf{Z-buffer (depth buffer)}, painter\textquotesingle s (depth sort), scan-line, BSP tree, back-face detection (N·V \textgreater{} 0 → back face), Warnock (area subdivision), ray casting.
\end{itemize}

\needspace{162pt}

\hypertarget{p2-03--7-deeper-dive--scoping-more-c-traps-java-essentials--graphics-traces}{%
\section{7. Deeper Dive --- Scoping, More C Traps, Java Essentials \& Graphics Traces}\label{p2-03--7-deeper-dive--scoping-more-c-traps-java-essentials--graphics-traces}}

\needspace{122pt}

\hypertarget{p2-03--71-static-vs-dynamic-scoping--classic-question}{%
\subsection{7.1 Static vs Dynamic Scoping --- Classic Question}\label{p2-03--71-static-vs-dynamic-scoping--classic-question}}

\begin{asciiblock}{67pt}
\begin{Verbatim}[fontsize=\fontsize{8.8}{10.7}\selectfont]
int x = 1;
void f()  { printf("%d", x); }
void g()  { int x = 2; f(); }
int main() { g(); }
\end{Verbatim}
\end{asciiblock}

\begin{itemize}
\tightlist
\item
  \textbf{Static (lexical) scoping}: f\textquotesingle s free variable x refers to the global x → prints \textbf{1} (C, Java, Pascal).
\item
  \textbf{Dynamic scoping}: x is looked up in the most recent active frame (g\textquotesingle s) → prints \textbf{2}.
\end{itemize}

\needspace{101pt}

\hypertarget{p2-03--72-recursion-trace}{%
\subsection{7.2 Recursion Trace}\label{p2-03--72-recursion-trace}}

\begin{asciiblock}{46pt}
\begin{Verbatim}[fontsize=\fontsize{8.8}{10.7}\selectfont]
int fun(int n) { if (n <= 1) return 1; return n * fun(n - 2); }
fun(7) = 7 × fun(5) = 7 × 5 × fun(3) = 7 × 5 × 3 × fun(1) = 105
\end{Verbatim}
\end{asciiblock}

\begin{asciiblock}{78pt}
\begin{Verbatim}[fontsize=\fontsize{8.8}{10.7}\selectfont]
Call stack (grows downward)       Returns (unwinding)
 fun(7)                            fun(1) → 1
   fun(5)                          fun(3) → 3 × 1  = 3
     fun(3)                        fun(5) → 5 × 3  = 15
       fun(1)                      fun(7) → 7 × 15 = 105
\end{Verbatim}
\end{asciiblock}

\needspace{362pt}

\hypertarget{p2-03--73-more-c-output-questions-with-reasons}{%
\subsection{7.3 More C Output Questions (with reasons)}\label{p2-03--73-more-c-output-questions-with-reasons}}

\begin{asciiblock}{307pt}
\begin{Verbatim}[fontsize=\fontsize{7.5}{9.1}\selectfont]
// 1. switch fall-through
int k = 2;
switch (k) { case 1: printf("A"); case 2: printf("B"); case 3: printf("C"); break; default: printf("D"); }
// Output: BC   (no break after case 2)

// 2. bitwise
printf("%d %d %d", 5 & 3, 5 | 3, 5 ^ 3);      // 1 7 6
printf("%d", ~5);                               // -6  (2's complement: ~x = -x - 1)

// 3. pointer to string
char *s = "UGCNET";
printf("%s", s + 3);                            // NET
printf("%c", *s + 1);                           // V   ('U' + 1)

// 4. post-increment in loop condition
int i = 0; while (i++ < 3); printf("%d", i);    // 4

// 5. integer division & casting
printf("%.2f", (float)7 / 2);                   // 3.50
printf("%d", 7 / 2 * 2);                        // 6

// 6. comma operator
int a = (1, 2, 3); printf("%d", a);             // 3

// 7. short-circuit
int p = 0, q = 5;
if (p && (q = 10)) {}  printf("%d", q);         // 5 (second operand never evaluated)

// 8. 2-D array pointer arithmetic
int m[2][3] = {{1,2,3},{4,5,6}};
printf("%d", *(*(m + 1) + 2));                  // 6  (m[1][2])
\end{Verbatim}
\end{asciiblock}

\needspace{135pt}

\hypertarget{p2-03--74-java-essentials}{%
\subsection{7.4 Java Essentials}\label{p2-03--74-java-essentials}}

\begingroup\small

\begin{longtable}[]{@{}
  >{\raggedright\arraybackslash}p{(\columnwidth - 8\tabcolsep) * \real{0.1237}}
  >{\raggedright\arraybackslash}p{(\columnwidth - 8\tabcolsep) * \real{0.1165}}
  >{\raggedright\arraybackslash}p{(\columnwidth - 8\tabcolsep) * \real{0.1367}}
  >{\raggedright\arraybackslash}p{(\columnwidth - 8\tabcolsep) * \real{0.1956}}
  >{\raggedright\arraybackslash}p{(\columnwidth - 8\tabcolsep) * \real{0.0645}}@{}}
\toprule\noalign{}
\begin{minipage}[b]{\linewidth}\raggedright
\textbf{Modifier}
\end{minipage} & \begin{minipage}[b]{\linewidth}\raggedright
\textbf{Same class}
\end{minipage} & \begin{minipage}[b]{\linewidth}\raggedright
\textbf{Same package}
\end{minipage} & \begin{minipage}[b]{\linewidth}\raggedright
\textbf{Subclass (other pkg)}
\end{minipage} & \begin{minipage}[b]{\linewidth}\raggedright
\textbf{World}
\end{minipage} \\
\midrule\noalign{}
\endhead
\bottomrule\noalign{}
\endlastfoot
private & ✔ & ✘ & ✘ & ✘ \\
default (none) & ✔ & ✔ & ✘ & ✘ \\
protected & ✔ & ✔ & ✔ & ✘ \\
public & ✔ & ✔ & ✔ & ✔ \\
\end{longtable}

\endgroup

\begingroup\small

\begin{longtable}[]{@{}
  >{\raggedright\arraybackslash}p{(\columnwidth - 2\tabcolsep) * \real{0.4090}}
  >{\raggedright\arraybackslash}p{(\columnwidth - 2\tabcolsep) * \real{0.4972}}@{}}
\toprule\noalign{}
\begin{minipage}[b]{\linewidth}\raggedright
\textbf{Abstract class}
\end{minipage} & \begin{minipage}[b]{\linewidth}\raggedright
\textbf{Interface}
\end{minipage} \\
\midrule\noalign{}
\endhead
\bottomrule\noalign{}
\endlastfoot
Can have constructors, state (fields), concrete methods & Constants + abstract methods (default/\allowbreak{}static methods since Java 8) \\
Single inheritance (\texttt{extends}) & A class can \texttt{implements} many interfaces \\
Use for "is-a" with shared code & Use for capability contracts \\
\end{longtable}

\endgroup

\begin{itemize}
\tightlist
\item
  \texttt{final} variable = constant; \texttt{final} method = cannot be overridden; \texttt{final} class = cannot be inherited (String).
\item
  \texttt{static} members belong to the class. \texttt{super} refers to the parent; \texttt{this} to the current object.
\item
  JVM components: \textbf{class loader → bytecode verifier → interpreter / JIT compiler}; runtime areas: heap, method area, stack, PC registers.
\item
  Exceptions: \textbf{checked} (IOException --- must be handled or declared) vs \textbf{unchecked} (RuntimeException, e.g. NullPointerException, ArithmeticException).
\item
  \texttt{String} is immutable; \texttt{StringBuffer} (synchronised) and \texttt{StringBuilder} (faster) are mutable.
\end{itemize}

\needspace{147pt}

\hypertarget{p2-03--75-graphics-traces}{%
\subsection{7.5 Graphics Traces}\label{p2-03--75-graphics-traces}}

\textbf{DDA}: (2, 3) → (8, 6); dx = 6, dy = 3, steps = 6, x-inc = 1, y-inc = 0.5.

\begin{asciiblock}{57pt}
\begin{Verbatim}[fontsize=\fontsize{8.8}{10.7}\selectfont]
x : 2   3    4   5    6   7    8
y : 3   3.5  4   4.5  5   5.5  6
plotted (round half up): (2,3) (3,4) (4,4) (5,5) (6,5) (7,6) (8,6)
\end{Verbatim}
\end{asciiblock}

\textbf{Mid-point circle, r = 10} (first octant, start (0, 10), p₀ = 1 − r = −9):

\begingroup\small

\begin{longtable}[]{@{}
  >{\raggedright\arraybackslash}p{(\columnwidth - 6\tabcolsep) * \real{0.0300}}
  >{\raggedright\arraybackslash}p{(\columnwidth - 6\tabcolsep) * \real{0.0385}}
  >{\raggedright\arraybackslash}p{(\columnwidth - 6\tabcolsep) * \real{0.1029}}
  >{\raggedright\arraybackslash}p{(\columnwidth - 6\tabcolsep) * \real{0.1070}}@{}}
\toprule\noalign{}
\begin{minipage}[b]{\linewidth}\raggedright
\textbf{k}
\end{minipage} & \begin{minipage}[b]{\linewidth}\raggedright
\textbf{pₖ}
\end{minipage} & \begin{minipage}[b]{\linewidth}\raggedright
\textbf{Next pixel}
\end{minipage} & \begin{minipage}[b]{\linewidth}\raggedright
\textbf{pₖ₊₁}
\end{minipage} \\
\midrule\noalign{}
\endhead
\bottomrule\noalign{}
\endlastfoot
0 & −9 & (1, 10) & −6 \\
1 & −6 & (2, 10) & −1 \\
2 & −1 & (3, 10) & 6 \\
3 & 6 & (4, 9) & −3 \\
4 & −3 & (5, 9) & 8 \\
5 & 8 & (6, 8) & 5 \\
6 & 5 & (7, 7) & stop (x ≥ y) \\
\end{longtable}

\endgroup

\textbf{Composite transformation}: rotate P(4, 2) by 90° anticlockwise about pivot (2, 2).

\begin{asciiblock}{57pt}
\begin{Verbatim}[fontsize=\fontsize{8.8}{10.7}\selectfont]
1. Translate by (−2, −2):  (2, 0)
2. Rotate 90°: (x cos90 − y sin90, x sin90 + y cos90) = (0, 2)
3. Translate by (+2, +2):  (2, 4)        → P' = (2, 4)
\end{Verbatim}
\end{asciiblock}

Scaling about fixed point (xf, yf): T(xf, yf) · S(sx, sy) · T(−xf, −yf) → x\textquotesingle{} = xf + (x − xf)·sx.

\textbf{Cohen--Sutherland}: window (0, 0)--(10, 10); line (−5, 5) → (15, 5). Codes 0001 and 0010; AND = 0000 → not trivially rejected; clip at x = 0 → (0, 5) and at x = 10 → (10, 5).

\textbf{Liang--Barsky}: same window; line (−5, 3) → (15, 9); Δx = 20, Δy = 6.

\begin{asciiblock}{86pt}
\begin{Verbatim}[fontsize=\fontsize{8.4}{10.2}\selectfont]
p₁ = −Δx = −20, q₁ = x₁ − xmin = −5   → r₁ = 0.25    (entering)
p₂ =  Δx =  20, q₂ = xmax − x₁ = 15   → r₂ = 0.75    (leaving)
p₃ = −Δy = −6,  q₃ = y₁ − ymin = 3    → r₃ = −0.5    (entering)
p₄ =  Δy =  6,  q₄ = ymax − y₁ = 7    → r₄ ≈ 1.17    (leaving)
u₁ = max(0, 0.25, −0.5) = 0.25 ;  u₂ = min(1, 0.75, 1.17) = 0.75
Clipped line: (−5 + 0.25·20, 3 + 0.25·6) = (0, 4.5)  to  (−5 + 0.75·20, 3 + 0.75·6) = (10, 7.5)
\end{Verbatim}
\end{asciiblock}

\needspace{165pt}

\hypertarget{p2-03--76-3-d-transformation-matrices-homogeneous-4--4}{%
\subsection{7.6 3-D Transformation Matrices (homogeneous 4 × 4)}\label{p2-03--76-3-d-transformation-matrices-homogeneous-4--4}}

\begin{asciiblock}{110pt}
\begin{Verbatim}[fontsize=\fontsize{8.8}{10.7}\selectfont]
Translation           Scaling              Rotation about z-axis
| 1 0 0 tx |          | sx 0  0  0 |       | cosθ −sinθ 0 0 |
| 0 1 0 ty |          | 0  sy 0  0 |       | sinθ  cosθ 0 0 |
| 0 0 1 tz |          | 0  0  sz 0 |       |  0     0   1 0 |
| 0 0 0 1  |          | 0  0  0  1 |       |  0     0   0 1 |
Rotation about x: y' = y cosθ − z sinθ, z' = y sinθ + z cosθ
Rotation about y: z' = z cosθ − x sinθ, x' = z sinθ + x cosθ
Perspective (centre of projection at origin, plane z = d): x' = x·d/z, y' = y·d/z
\end{Verbatim}
\end{asciiblock}

\needspace{135pt}

\hypertarget{p2-03--previous-year-questions-pyq-pattern}{%
\section{Previous Year Questions (PYQ pattern)}\label{p2-03--previous-year-questions-pyq-pattern}}

\textbf{Language concepts}

\begin{enumerate}
\def\labelenumi{\arabic{enumi}.}
\tightlist
\item
  First object-oriented language: \textbf{Simula 67}
\item
  Which uses call by name? \textbf{ALGOL 60}
\item
  Binding of a variable to its type in C occurs at: \textbf{compile time}
\item
  Dynamic scoping resolves a non-local name using: \textbf{the calling sequence (most recent active binding)}
\item
  Prolog is a: \textbf{logic programming language}
\item
  LISP is primarily used for: \textbf{list processing / AI (functional)}
\end{enumerate}

\textbf{C}

\begin{enumerate}
\def\labelenumi{\arabic{enumi}.}
\setcounter{enumi}{6}
\tightlist
\item
  \texttt{int\ a\ =\ 5;\ a\ =\ a++\ +\ ++a;} → \textbf{undefined behaviour} (option often given as 12 or 13; correct per standard: UB)
\item
  Output of \texttt{printf("\%d",\ sizeof(\textquotesingle{}A\textquotesingle{}))} in C: \textbf{4} (char constant is int in C; 1 in C++)
\item
  \texttt{\#define\ MUL(a,b)\ a*b}; \texttt{MUL(2+3,\ 4)} → 2+3*4 = \textbf{14}
\item
  Size of union \texttt{\{int\ a;\ char\ b{[}10{]};\}} (int=4): \textbf{12} with alignment (largest member 10, padded to multiple of 4)
\item
  Which storage class retains value between function calls? \textbf{static}
\item
  \texttt{char\ *p\ =\ "hello";\ printf("\%c",\ *(p+1));} → \textbf{e}
\item
  \texttt{int\ (*p){[}5{]};} declares: \textbf{pointer to an array of 5 integers}
\item
  Which function allocates zero-initialised memory? \textbf{calloc}
\end{enumerate}

\textbf{OOP/C++}

\begin{enumerate}
\def\labelenumi{\arabic{enumi}.}
\setcounter{enumi}{14}
\tightlist
\item
  Run-time polymorphism is achieved via: \textbf{virtual functions}
\item
  A class with at least one pure virtual function is: \textbf{abstract class}
\item
  Which operator cannot be overloaded? \textbf{:: (scope resolution)} (also \texttt{.}, \texttt{?:}, \texttt{sizeof})
\item
  The copy constructor takes argument by: \textbf{reference} (by value would cause infinite recursion)
\item
  Diamond problem is resolved by: \textbf{virtual base class}
\item
  Templates support: \textbf{generic programming (compile-time polymorphism)}
\item
  Order of destructor calls: \textbf{reverse of constructor calls}
\item
  Java does not support multiple inheritance of classes; achieved via: \textbf{interfaces}
\end{enumerate}

\textbf{Web}

\begin{enumerate}
\def\labelenumi{\arabic{enumi}.}
\setcounter{enumi}{22}
\tightlist
\item
  Applet life cycle method called first: \textbf{init()}
\item
  Servlet method handling each request: \textbf{service()}
\item
  XML document conforming to DTD is called: \textbf{valid}; following syntax rules: \textbf{well-formed}
\item
  DHTML combines: \textbf{HTML, CSS, JavaScript, DOM}
\end{enumerate}

\textbf{Graphics}

\begin{enumerate}
\def\labelenumi{\arabic{enumi}.}
\setcounter{enumi}{26}
\tightlist
\item
  Frame buffer for 640×480 with 8 bits per pixel: \textbf{300 KB} (307200 bytes)
\item
  Bresenham\textquotesingle s line algorithm uses: \textbf{only integer arithmetic}
\item
  Initial decision parameter for mid-point circle with r = 10: \textbf{1 − 10 = −9}
\item
  Circle uses \_\_-way symmetry, ellipse \_\_-way: \textbf{8, 4}
\item
  Composite transformation for rotation about point P: \textbf{T(P)·R(θ)·T(−P)}
\item
  Region code of a point above-left of window (TBRL): \textbf{1001}
\item
  Line endpoints codes 0101 and 0110: AND = 0100 ≠ 0 → \textbf{trivially rejected}
\item
  Polygon clipping algorithm: \textbf{Sutherland--Hodgman}
\item
  Shading that interpolates normal vectors: \textbf{Phong}
\item
  Z-buffer algorithm is a: \textbf{image-space hidden surface method}
\item
  Bezier curve with 4 control points has degree: \textbf{3}
\item
  Which curve has local control? \textbf{B-spline}
\item
  Projection in which parallel lines converge: \textbf{Perspective}
\item
  Isometric projection is a type of: \textbf{axonometric orthographic projection}
\item
  Scaling matrix with sx = sy = −1 is equivalent to: \textbf{reflection about origin (rotation by 180°)}
\item
  Translating (2,3) by (4,−1): \textbf{(6,2)}
\end{enumerate}

\textbf{More practice questions}

\begin{enumerate}
\def\labelenumi{\arabic{enumi}.}
\setcounter{enumi}{42}
\tightlist
\item
  Under dynamic scoping, the code in §7.1 prints: \textbf{2}
\item
  Output of \texttt{printf("\%d",\ \textasciitilde{}0);} in C: \textbf{−1}
\item
  Output of \texttt{printf("\%d",\ 10\ \textgreater{}\textgreater{}\ 1\ \textless{}\textless{}\ 2);}: (10 \textgreater\textgreater{} 1) \textless\textless{} 2 = \textbf{20}
\item
  In Java, a member visible only within its package has: \textbf{default (package-private) access}
\item
  A Java class that cannot be subclassed is declared: \textbf{final}
\item
  NullPointerException is a: \textbf{unchecked (runtime) exception}
\item
  Number of steps in DDA for line (1, 1) to (9, 4): \textbf{8}
\item
  Rotating point (1, 0) by 90° anticlockwise about origin gives: \textbf{(0, 1)}
\item
  Reflection of (3, −2) about the line y = x: \textbf{(−2, 3)}
\item
  In Liang--Barsky, pₖ \textless{} 0 means the line is: \textbf{entering} the clip boundary
\item
  Shear in x with shx = 2 applied to (1, 3): \textbf{(7, 3)}
\item
  Homogeneous coordinates for 3-D points use: \textbf{4-element vectors / 4 × 4 matrices}
\end{enumerate}

\begin{revbox}

\begin{itemize}
\tightlist
\item
  Simula 67 first OO · ALGOL call-by-name · C static scope
\item
  \texttt{::}, \texttt{.}, \texttt{?:}, \texttt{sizeof} can\textquotesingle t be overloaded · Pure virtual → abstract
\item
  Bresenham p0 = 2dy − dx · Mid-point circle p0 = 1 − r
\item
  TBRL codes; AND ≠ 0 reject; both 0 accept
\item
  Gouraud = intensity interpolation · Phong = normal interpolation
\item
  Bezier global control, passes endpoints · B-spline local control
\end{itemize}

\end{revbox}
